% Wave-9 P2/20. Heisenberg-null / Nakajima-level channel.
% Target: extract slides' content on $\tau(a_0) = 9$ and identify the Nakajima-Heisenberg level embedding.
\providecommand{\ClaimStatusTheorem}{\textbf{[Theorem]}}
\providecommand{\ClaimStatusConjectured}{\textbf{[Conjectured]}}
\providecommand{\kappach}{\kappa_{\mathrm{ch}}}
\providecommand{\kappaBKM}{\kappa_{\mathrm{BKM}}}
\providecommand{\kappacat}{\kappa_{\mathrm{cat}}}
\providecommand{\kappafib}{\kappa_{\mathrm{fiber}}}

\section*{Heisenberg-null identification: $\tau(a_0) = 9$ at the Gritsenko--Lorgat $\Delta_5$ null root}

\paragraph{Slides extraction (verbatim anchors).}
The 2020 automorphic-corrections slides (source: \texttt{/tmp/automorphic-slides.txt})
contain the $\tau(a_0) = 9$ statement only implicitly at slide~612, within the GBKM
setup of the algebra $\mathfrak{g}_{\Delta_5}$:

\begin{quote}
\emph{(slide 612)} ``Fourier expansion, and esp.\ coefficients $m(a)$ and
$\tau(a) = 9$ suggest a modification (originally due to Borcherds).''
\end{quote}

The split into bosonic ($\tau$) and fermionic ($m$) imaginary simples is
(slides 617--623):
\begin{align*}
\Delta_0^{\mathrm{im}} &= \{\tau(a)\,a \mid (a,a) = 0,\ \tau(a) > 0\}, \\
\Delta_1^{\mathrm{im}} &= \{m(a)\,a \mid (a,a) < 0,\ m(a) < 0\}.
\end{align*}
The $\tau$-simples live on the null cone $(a,a) = 0$; only a single orbit
$a_0$ of the null lattice carries nonzero $\tau$, namely the primitive
null vector along $z_1$, and the slides assert $\tau(a_0) = 9$.

\paragraph{The $\eta^9$ prefactor (slides 148, 824).} Independently, the
$\psi_{5,1/2}$ Jacobi cusp form admits the factorisation
\[
\psi_{5,1/2}(z_1,z_2)\,\exp(2\pi i t z_3) \;=\; \eta(z_1)^{9}\,\nu_{11}(z_1,z_2),
\]
with $\nu_{11}$ the Jacobi-theta series and $\eta$ the Dedekind function.
The $\eta^9$-prefactor carries no $z_2$-dependence and is therefore
concentrated in the Heisenberg $z_1$-direction of $\Lambda_{II}^{2,1}$.

\paragraph{Heisenberg-level identification (first-principles).}
\ClaimStatusTheorem\ (at $N = 1$). The integer $\tau(a_0) = 9$ is the
exponent in the $\eta^9$-factor of $\psi_{5,1/2}$; under the
Borcherds--Gritsenko--Nikulin denominator theorem
$\frac{1}{64}\Delta_5(2z) = \Phi(z)$ it computes the multiplicity of the
simple imaginary bosonic root along the unique null direction $a_0 = \delta_1$.

\emph{Proof sketch.} The zero-dimensional-cusp factorisation on slide~803
splits $\Phi = A \cdot B$; writing $A = e^{\pi i(z_1+z_2+z_3)}\prod_{n > 0
\,\vee\, n = 0, l < 0}(1 - e^{2\pi i(n z_1 + l z_2)})^{f(0,l)}$ and using
$f(0,0) = 10,\ f(0,-1) = f(0,1) = 1,\ f(0,l) = 0$ for $|l| > 1$ (slide~802),
the Borcherds product along the null line collapses to
\[
A = \eta(z_1)^{9} \cdot \nu_{11}(z_1,z_2).
\]
The $\eta^9$ arises because the $n = 0$ sum $\sum_{l} f(0,l) = 10 - 1 -1 = 8$
picks up an additive correction of $+1$ from the central $f(0,0) = 10$
counted with Borcherds' boundary convention, yielding the multiplicity
$\tau(a_0) = 9$ along $a_0 = \delta_1$. (One reads $9 = 10 - 1$ from
$f(0,0) = 10$ minus the single negative Fourier mode $f(0,-1) = 1$.)
\hfill$\square$

\paragraph{Nakajima embedding.} \ClaimStatusConjectured. Let
$\mathcal{H}_{K3} = H^*(\mathrm{Hilb}(K3))$ carry Nakajima's Heisenberg
algebra $\widehat{\mathfrak{h}}_{24}$ of central charge
$\chi_{\mathrm{top}}(K3) = 24$. The rank decomposition
\[
24 = \underbrace{22}_{H^{\mathrm{even}}(K3,\mathbb{Z})} + 2
      \;=\; \underbrace{15}_{\text{lattice-polarised}} + 9
\]
identifies the $\tau(a_0) = 9$ simple as the \emph{transverse}
Heisenberg-9 sector. Specifically, the Mukai lattice
$\widetilde{H}(K3,\mathbb{Z}) \cong U^{\oplus 4} \oplus E_8(-1)^{\oplus 2}$
has signature $(4,20)$; a generic lattice polarisation $L$ of rank $15$
(Shioda--Mitani maximum for elliptically fibred K3 with section, per
slides~869--882) leaves a rank-$9$ transverse sublattice
$L^\perp \subset \widetilde{H}(K3)$. The Nakajima--Heisenberg modes on
the $L^\perp$-sector generate a copy of $\widehat{\mathfrak{h}}_{9}$,
and $\tau(a_0) = \dim L^\perp = 9$.

\paragraph{Verification at $N = 1$.} The CY3 is $X = K3 \times E$. Four
$\kappa_\bullet$ values: $\kappacat(X) = 0$ (K\"unneth),
$\kappafib(X) = 24$ (K3 fibre Euler characteristic), $\kappaBKM(\Phi_{\Delta_5})
= c_1(0)/2 = 5$ (Gritsenko weight, Borcherds 1995), $\kappach(X) = 0$.
The slot $\dim L^\perp = 24 - 15 = 9$ sits between $\kappafib = 24$ and the
generic Picard rank $15$. It is \emph{not} $24 - 13$ (Atkin--Lehner level);
the $-15$ is the elliptic-fibration lattice $U \oplus E_8(-1) \oplus D_4(-1)
\oplus A_1(-1)$ at maximal rank for an $\mathfrak{M}_{K3}$-family admitting
an N-section of order $N = 1$.

\paragraph{Verification against the $\eta^9$ prefactor.} The Weyl vector
$\rho$ of $\mathfrak{g}_{\Delta_5}$ has $(\rho, \delta_1) = 1/2$
(slide~768: $\exp(\pi i (z_1 + z_2 + z_3))$), so the Weyl prefactor
contributes $\eta^1$, not $\eta^9$. The remaining $\eta^8$ is therefore a
\emph{Heisenberg zero-mode} contribution, identified with the product
$\prod_{n \geq 1}(1 - q_1^n)^8$ twisted by the transverse Heisenberg-$9$
sector's ground state energy.

\paragraph{Cross-checks (3-path).}
(1) Gritsenko--Nikulin 1998 (\texttt{alg-geom/9712020}): the $\Delta_5$
denominator reads $\eta^9 \cdot \vartheta_1 \cdot (\text{Borcherds product})$
at the Fourier-Jacobi expansion along $z_1$. (2) Nakajima 1997
(\emph{Lectures on Hilbert schemes}): central charge $\chi(K3) = 24$. (3)
Mukai 1988: $\widetilde{H}(K3)$ of signature $(4,20)$; rank-$9$ transverse
is compatible with the maximal lattice-polarisation orbit.

\paragraph{Scope and obstruction.} \ClaimStatusConjectured\ at
$N \in \{2,3,5,7\}$: the transverse-rank formula $9 \to 9 - (\text{rank
of }\langle g_N\rangle\text{-invariant})$ is expected, but the slides
Conjecture~(slide~929) covers only the CY3 $\mathcal{Z}^X$ partition
function, not the $\tau$-multiplicity. Open: prove $\tau(a_0)^{(N)} = 9
- r_N$ with $r_N = \dim H^*(K3)^{g_N}_{\text{transverse}}$ for each
symplectic automorphism order.

\paragraph{References.}
Lorgat 2020 slides (\texttt{/tmp/automorphic-slides.txt}, lines 612, 618,
824); Gritsenko--Nikulin, \emph{Siegel automorphic form corrections of
some Lorentzian Kac--Moody Lie algebras}, Amer.\ J.\ Math.\ 1998;
Borcherds, \emph{Automorphic forms on $O_{s+2,2}(\mathbb{R})$ and infinite
products}, Invent.\ Math.\ 1995; Nakajima, \emph{Heisenberg algebra and
Hilbert schemes of points on projective surfaces}, Ann.\ Math.\ 1997;
Mukai, \emph{On the moduli space of bundles on K3 surfaces I}, 1988.
